\documentclass[12pt,a4paper]{article}
\usepackage[utf8]{vietnam}
\usepackage{amsmath,amssymb}
\usepackage{geometry}
\geometry{
  top=12mm,
  bottom=12mm,
  left=14mm,
  right=14mm,
  footskip=6mm,
  headheight=14pt
}

\usepackage{tikz}
\usetikzlibrary{arrows.meta,calc,shapes.geometric}
\usepackage[most]{tcolorbox}
\usepackage{fontawesome}
\usepackage{fancyhdr}
\usepackage{enumitem}

\pagestyle{fancy}
\fancyhf{}
\fancyhead[L]{\footnotesize\color{teal!80!black}\faBook\ \textbf{SỔ TAY KIẾN THỨC TOÁN 3}}
\fancyhead[R]{\footnotesize\color{teal!80!black}\textbf{Lớp Toán Thầy Lĩnh -- loptoanthaylinh.vn}}
\fancyfoot[C]{\footnotesize\color{gray}-- Sổ tay kiến thức: Bài 16. Điểm ở giữa, trung điểm của đoạn thẳng --}
\renewcommand{\headrulewidth}{0.4pt}
\renewcommand{\footrulewidth}{0pt}

\setlength{\parindent}{0pt}
\setlength{\parskip}{0pt}

\begin{document}

% --- BANNER BÀI HỌC ---
\begin{tcolorbox}[
  enhanced,
  colback=teal!8!white,
  colframe=teal!70!black,
  arc=3mm, boxrule=1.2pt,
  left=3mm, right=3mm, top=2mm, bottom=2mm,
  halign=center
]
  {\LARGE\bfseries\color{teal!90!black}\faBookmark\ BÀI 16: ĐIỂM Ở GIỮA, TRUNG ĐIỂM CỦA ĐOẠN THẲNG}\\[1mm]
  {\itshape\footnotesize\color{black!75} Phân biệt điểm ở giữa và trung điểm, kĩ thuật chia đôi đoạn thẳng trên 1 trang A4}
\end{tcolorbox}

\vspace{-2mm}

\noindent
\begin{minipage}[t]{0.49\textwidth}\vspace{0pt}
\begin{tcolorbox}[
  enhanced,
  colback=blue!4!white,
  colframe=blue!75!black,
  arc=2mm, boxrule=0.8pt,
  left=2mm, right=2mm, top=1.8mm, bottom=1.8mm,
  title={\footnotesize\bfseries\color{white}\faTags\ 1. Điểm ở giữa}
]
{\footnotesize
\begin{itemize}\setlength{\itemsep}{1pt}\setlength{\leftskip}{-2.5mm}
  \item \textbf{Điều kiện:} Ba điểm $A, B, C$ cùng nằm trên một đường thẳng (\textbf{thẳng hàng}) theo thứ tự đó $\implies$ $B$ là \textbf{điểm ở giữa} hai điểm $A$ và $C$.
  \item Điểm ở giữa không nhất thiết phải chia đôi đoạn thẳng (chỉ cần nằm giữa hai điểm mút).
  \item Nếu $A, B, C$ không thẳng hàng thì $B$ \textbf{không} được gọi là điểm ở giữa.
\end{itemize}
}
\end{tcolorbox}
\end{minipage}
\hfill
\begin{minipage}[t]{0.49\textwidth}\vspace{0pt}
\begin{tcolorbox}[
  enhanced,
  colback=teal!4!white,
  colframe=teal!75!black,
  arc=2mm, boxrule=0.8pt,
  left=2mm, right=2mm, top=1.8mm, bottom=1.8mm,
  title={\footnotesize\bfseries\color{white}\faBalanceScale\ 2. Trung điểm của đoạn thẳng}
]
{\footnotesize
\begin{itemize}\setlength{\itemsep}{1pt}\setlength{\leftskip}{-2.5mm}
  \item Điểm $M$ là \textbf{trung điểm} của đoạn thẳng $AB$ khi thỏa mãn \textbf{đủ 2 điều kiện}:
  \begin{enumerate}\setlength{\itemsep}{0.5pt}\setlength{\leftskip}{-2.5mm}
    \item $M$ là \textbf{điểm ở giữa} hai điểm $A$ và $B$ ($A, M, B$ thẳng hàng).
    \item Độ dài hai đoạn bằng nhau: $\mathbf{AM = MB}$.
  \end{enumerate}
  \item \textbf{Công thức độ dài:}
  $$AM = MB = \dfrac{AB}{2} \iff AB = 2 \times AM$$
\end{itemize}
}
\end{tcolorbox}
\end{minipage}

\vspace{-2.5mm}
% --- KHỐI 3: SƠ ĐỒ TƯ DUY TIKZ ---
\noindent
\begin{tcolorbox}[
  enhanced,
  colback=orange!4!white,
  colframe=orange!80!black,
  arc=2mm, boxrule=0.8pt,
  left=2mm, right=2mm, top=1.5mm, bottom=1.5mm,
  title={\footnotesize\bfseries\color{white}\faSitemap\ SƠ ĐỒ TƯ DUY: HÌNH HỌC TRỰC QUAN TRUNG ĐIỂM ĐOẠN THẲNG}
]
\centering
\begin{tikzpicture}[
  dot/.style={circle, fill=blue!80!black, inner sep=1.5pt}
]
  \draw[thick, blue!80!black] (0, 0) -- (8, 0);
  \node[dot] (A) at (0, 0) {};
  \node[dot, fill=red!80!black] (M) at (4, 0) {};
  \node[dot] (B) at (8, 0) {};

  \node[below, font=\scriptsize\bfseries] at (A) {$A$};
  \node[below, font=\scriptsize\bfseries, red!80!black] at (M) {$M$ (Trung điểm)};
  \node[below, font=\scriptsize\bfseries] at (B) {$B$};

  % Tick marks
  \draw[thick, teal!80!black] (1.9, -0.15) -- (1.9, 0.15);
  \draw[thick, teal!80!black] (2.1, -0.15) -- (2.1, 0.15);
  \draw[thick, teal!80!black] (5.9, -0.15) -- (5.9, 0.15);
  \draw[thick, teal!80!black] (6.1, -0.15) -- (6.1, 0.15);

  \node[above, font=\scriptsize\bfseries, teal!80!black] at (2, 0.2) {$AM = 4\text{ cm}$};
  \node[above, font=\scriptsize\bfseries, teal!80!black] at (6, 0.2) {$MB = 4\text{ cm}$};
  \draw[<->, purple!80!black, thick] (0, -0.6) -- (8, -0.6) node[midway, fill=orange!4!white, font=\scriptsize\bfseries] {Đoạn thẳng $AB = 8\text{ cm}$};
\end{tikzpicture}
\end{tcolorbox}

\vspace{-2.5mm}
% --- KHỐI 4: VÍ DỤ & MẸO TRÁNH BẪY ---
\noindent
\begin{minipage}[t]{0.49\textwidth}\vspace{0pt}
\begin{tcolorbox}[
  enhanced,
  colback=green!4!white,
  colframe=green!60!black,
  arc=2mm, boxrule=0.8pt,
  left=2mm, right=2mm, top=1.8mm, bottom=1.8mm,
  title={\footnotesize\bfseries\color{white}\faCheckCircle\ Ví dụ thực hành chuẩn mực}
]
{\footnotesize
\textbf{1. Tính độ dài:}\\
Đoạn thẳng $AB$ dài $10\text{ cm}$, điểm $M$ là trung điểm của đoạn thẳng $AB$. Tính độ dài đoạn thẳng $AM$.\\
\textbf{Bài giải:}\\
Độ dài đoạn thẳng $AM$ là:
$$10 : 2 = \mathbf{5}\text{ (cm)}$$
\textbf{2. Tìm độ dài cả đoạn:}\\
Biết $I$ là trung điểm đoạn thẳng $CD$ và $CI = 6\text{ cm}$. Độ dài đoạn thẳng $CD$ là: $6 \times 2 = \mathbf{12}\text{ cm}$.
}
\end{tcolorbox}
\end{minipage}
\hfill
\begin{minipage}[t]{0.49\textwidth}\vspace{0pt}
\begin{tcolorbox}[
  enhanced,
  colback=red!4!white,
  colframe=red!70!black,
  arc=2mm, boxrule=0.8pt,
  left=2mm, right=2mm, top=1.8mm, bottom=1.8mm,
  title={\footnotesize\bfseries\color{white}\faExclamationTriangle\ Mẹo tránh bẫy trắc nghiệm}
]
{\footnotesize
\begin{itemize}\setlength{\itemsep}{1.5pt}\setlength{\leftskip}{-2.5mm}
  \item \textbf{Phân biệt 'Điểm ở giữa' và 'Trung điểm':}\\
  Mọi trung điểm đều là điểm ở giữa, nhưng điểm ở giữa chưa chắc đã là trung điểm (vì hai đoạn có thể dài ngắn khác nhau).
  \item \textbf{Bẫy thẳng hàng:}\\
  Nếu $MA = MB$ nhưng 3 điểm $A, M, B$ tạo thành một hình tam giác (không thẳng hàng) thì $M$ \textbf{không phải} là trung điểm của $AB$.
\end{itemize}
}
\end{tcolorbox}
\end{minipage}


\end{document}
